\section{Team dynamics}

Reporting on performance can be a thorny task: 
A successful assessment will allow the reader to make conclusions where and why
and what performance improvements would be beneficial to the code. When we write the report, it
is therefore important to stick to a factual, constructive language: Scientists
tend to be defensive when it comes to their code and often will start to explain why their particular code
has to behave this and that way.



These quickly become non-constructive conversations.
They are tiring and can lead to situations where we don't get the appreciation
for the work invested in the performance assessment.
Potential lessons learned are not taken on board. 
A performance assessment is simply a report how good an algorithm exploits a
machine for one particular setup, and this has to come across in any statement
on the code.
There should be little room to challenge the assessment, and this is best
achieved with a very factual, technical language.
